
\documentclass{article}
\usepackage{colortbl}
\usepackage{makecell}
\usepackage{multirow}
\usepackage{supertabular}

\begin{document}

\newcounter{utterance}

\twocolumn

{ \footnotesize  \setcounter{utterance}{1}
\setlength{\tabcolsep}{0pt}
\begin{supertabular}{c@{$\;$}|p{.15\linewidth}@{}p{.15\linewidth}p{.15\linewidth}p{.15\linewidth}p{.15\linewidth}p{.15\linewidth}}

    \# & $\;$A & \multicolumn{4}{c}{Game Master} & $\;\:$B\\
    \hline 

    \theutterance \stepcounter{utterance}  

    & & \multicolumn{4}{p{0.6\linewidth}}{\cellcolor[rgb]{0.9,0.9,0.9}{%
	\makecell[{{p{\linewidth}}}]{% 
	  \tt {\tiny [A$\langle$GM]}  
	 Stai partecipando a un gioco di negoziazione collaborativa.\\ \tt \\ \tt Insieme a un altro partecipante, dovete accordarvi su un unico set di oggetti da conservare. Ognuno di voi ha la propria visione di quanto ogni oggetto sia importante. Non sai come l'altro partecipante valuta gli oggetti. Inoltre, ti viene fornito lo sforzo che ogni oggetto richiede.\\ \tt Potete accordarvi su un set solo se lo sforzo totale degli oggetti selezionati non supera un limite condiviso:\\ \tt \\ \tt LIMITE = 2307\\ \tt \\ \tt Ecco i valori di sforzo dei singoli oggetti:\\ \tt \\ \tt Sforzo oggetto = {"A21": 147, "B93": 698, "A11": 588, "B33": 43, "C08": 356, "B39": 80, "B96": 95, "C85": 743, "A82": 106, "B31": 308, "B20": 325, "A60": 255, "B09": 276, "C81": 543, "B21": 51}\\ \tt \\ \tt Ecco la tua visione personale sull'importanza di ogni oggetto:\\ \tt \\ \tt Valori di importanza dell'oggetto = {"A21": 138, "B93": 583, "A11": 65, "B33": 262, "C08": 121, "B39": 508, "B96": 461, "C85": 484, "A82": 668, "B31": 389, "B20": 215, "A60": 97, "B09": 500, "C81": 30, "B21": 400}\\ \tt \\ \tt Obiettivo:\\ \tt \\ \tt Il tuo obiettivo è negoziare un set condiviso di oggetti che ti avvantaggi il più possibile (cioè, massimizzi l'importanza totale per TE), rimanendo entro il LIMITEE. Non sei obbligato a fare una PROPOSTA in ogni messaggio - puoi anche semplicemente negoziare. Tutte le tattiche sono permesse!\\ \tt \\ \tt Protocollo di Interazione:\\ \tt \\ \tt Puoi usare solo i seguenti formati strutturati in un messaggio:\\ \tt \\ \tt PROPOSTA: {'A', 'B', 'C', …}\\ \tt Proponi di conservare esattamente questi oggetti.\\ \tt RIFIUTO: {'A', 'B', 'C', …}\\ \tt Rifiuta esplicitamente la proposta dell'avversario.\\ \tt ARGOMENTAZIONE: {'...'}\\ \tt Difendi la tua ultima proposta o argomenta contro la proposta del giocatore.\\ \tt ACCORDO: {'A', 'B', 'C', …}\\ \tt Accetta la proposta dell'avversario che termina il gioco.\\ \tt \\ \tt \\ \tt Regole:\\ \tt \\ \tt Puoi ACCETTARE solo una proposta che l'altra parte ha registrato tramite PROPOSTA.\\ \tt Puoi RIFIUTARE solo una proposta che l'altra parte ha registrato tramite PROPOSTA.\\ \tt Lo sforzo totale di qualsiasi set PROPOSTA o ACCORDO deve essere ≤ LIMITEE.\\ \tt NON rivelare i tuoi punteggi di importanza nascosti.\\ \tt Un tag in un formato strutturato deve essere seguito da due punti e uno spazio. L'argomento deve essere un set python contenente 0 o più stringhe.\\ \tt Quindi, deve essere della forma TAG: {...}\\ \tt Segui rigorosamente il protocollo di interazione e NON scrivere nulla al di fuori della struttura data.\\ \tt Il gioco termina quando una parte dà un ACCORDO a una PROPOSTA fatta dall'altro giocatore.\\ \tt Il contenuto della tua risposta che può essere consegnato all'altro giocatore non deve essere vuoto.\\ \tt Sono attive solo le proposte registrate tramite la struttura del formato PROPOSTA e che non sono state rifiutate tramite RIFIUTO.\\ \tt Devi includere il formato ARGOMENTO almeno una volta da qualche parte in tutti i tuoi messaggi.\\ \tt  
	  } 
	   } 
	   } 
	 & \\ 
 

    \theutterance \stepcounter{utterance}  

    & \multicolumn{4}{p{0.6\linewidth}}{\cellcolor[rgb]{0.8,1,0.9}{%
	\makecell[{{p{\linewidth}}}]{% 
	  \tt {\tiny [A$\rangle$GM]}  
	 PROPOSTA: {'A21', 'B93', 'A11', 'B33', 'C08', 'B39', 'B96', 'C85', 'A82', 'B31', 'B20', 'A60', 'B09', 'C81', 'B21'}\\ \tt \\ \tt ARGOMENTAZIONE: {'Aspetto la tua risposta, il mio obiettivo è quello di massimizzare l'importanza totale, Cercherò di ottimizzare la scelta degli oggetti per non superare il LIMITE di 2307'} 
	  } 
	   } 
	   } 
	 & & \\ 
 

    \theutterance \stepcounter{utterance}  

    & & \multicolumn{4}{p{0.6\linewidth}}{\cellcolor[rgb]{0.9,0.9,0.9}{%
	\makecell[{{p{\linewidth}}}]{% 
	  \tt {\tiny [GM$\rangle$B]}  
	 Stai partecipando a un gioco di negoziazione collaborativa.\\ \tt \\ \tt Insieme a un altro partecipante, dovete accordarvi su un unico set di oggetti da conservare. Ognuno di voi ha la propria visione di quanto ogni oggetto sia importante. Non sai come l'altro partecipante valuta gli oggetti. Inoltre, ti viene fornito lo sforzo che ogni oggetto richiede.\\ \tt Potete accordarvi su un set solo se lo sforzo totale degli oggetti selezionati non supera un limite condiviso:\\ \tt \\ \tt LIMITE = 2307\\ \tt \\ \tt Ecco i valori di sforzo dei singoli oggetti:\\ \tt \\ \tt Sforzo oggetto = {"A21": 147, "B93": 698, "A11": 588, "B33": 43, "C08": 356, "B39": 80, "B96": 95, "C85": 743, "A82": 106, "B31": 308, "B20": 325, "A60": 255, "B09": 276, "C81": 543, "B21": 51}\\ \tt \\ \tt Ecco la tua visione personale sull'importanza di ogni oggetto:\\ \tt \\ \tt Valori di importanza dell'oggetto = {"A82": 30, "B93": 65, "B39": 97, "B09": 121, "C85": 138, "B96": 215, "B21": 262, "B31": 389, "B33": 400, "B20": 461, "A21": 484, "C08": 500, "A60": 508, "A11": 583, "C81": 668}\\ \tt \\ \tt Obiettivo:\\ \tt \\ \tt Il tuo obiettivo è negoziare un set condiviso di oggetti che ti avvantaggi il più possibile (cioè, massimizzi l'importanza totale per TE), rimanendo entro il LIMITEE. Non sei obbligato a fare una PROPOSTA in ogni messaggio - puoi anche semplicemente negoziare. Tutte le tattiche sono permesse!\\ \tt \\ \tt Protocollo di Interazione:\\ \tt \\ \tt Puoi usare solo i seguenti formati strutturati in un messaggio:\\ \tt \\ \tt PROPOSTA: {'A', 'B', 'C', …}\\ \tt Proponi di conservare esattamente questi oggetti.\\ \tt RIFIUTO: {'A', 'B', 'C', …}\\ \tt Rifiuta esplicitamente la proposta dell'avversario.\\ \tt ARGOMENTAZIONE: {'...'}\\ \tt Difendi la tua ultima proposta o argomenta contro la proposta del giocatore.\\ \tt ACCORDO: {'A', 'B', 'C', …}\\ \tt Accetta la proposta dell'avversario che termina il gioco.\\ \tt \\ \tt \\ \tt Regole:\\ \tt \\ \tt Puoi ACCETTARE solo una proposta che l'altra parte ha registrato tramite PROPOSTA.\\ \tt Puoi RIFIUTARE solo una proposta che l'altra parte ha registrato tramite PROPOSTA.\\ \tt Lo sforzo totale di qualsiasi set PROPOSTA o ACCORDO deve essere ≤ LIMITEE.\\ \tt NON rivelare i tuoi punteggi di importanza nascosti.\\ \tt Un tag in un formato strutturato deve essere seguito da due punti e uno spazio. L'argomento deve essere un set python contenente 0 o più stringhe.\\ \tt Quindi, deve essere della forma TAG: {...}\\ \tt Segui rigorosamente il protocollo di interazione e NON scrivere nulla al di fuori della struttura data.\\ \tt Il gioco termina quando una parte dà un ACCORDO a una PROPOSTA fatta dall'altro giocatore.\\ \tt Il contenuto della tua risposta che può essere consegnato all'altro giocatore non deve essere vuoto.\\ \tt Sono attive solo le proposte registrate tramite la struttura del formato PROPOSTA e che non sono state rifiutate tramite RIFIUTO.\\ \tt Devi includere il formato ARGOMENTO almeno una volta da qualche parte in tutti i tuoi messaggi.\\ \tt \\ \tt \\ \tt Ora riceverai il primo messaggio dell'altro giocatore.\\ \tt \\ \tt PROPOSTA: {'A21', 'B93', 'A11', 'B33', 'C08', 'B39', 'B96', 'C85', 'A82', 'B31', 'B20', 'A60', 'B09', 'C81', 'B21'}\\ \tt \\ \tt ARGOMENTAZIONE: {'Aspetto la tua risposta, il mio obiettivo è quello di massimizzare l'importanza totale, Cercherò di ottimizzare la scelta degli oggetti per non superare il LIMITE di 2307'} 
	  } 
	   } 
	   } 
	 & \\ 
 

    \theutterance \stepcounter{utterance}  

    & & & \multicolumn{4}{p{0.6\linewidth}}{\cellcolor[rgb]{1,0.85,0.72}{%
	\makecell[{{p{\linewidth}}}]{% 
	  \tt {\tiny [GM$\langle$B]}  
	 RIFIUTO: {'A21', 'B93', 'A11', 'B33', 'C08', 'B39', 'B96', 'C85', 'A82', 'B31', 'B20', 'A60', 'B09', 'C81', 'B21'}\\ \tt \\ \tt ARGOMENTAZIONE: {'La proposta iniziale supera il LIMITE di 2307, quindi non è valida. Sono disposto a negoziare per trovare un set di oggetti che massimizzi l'importanza totale senza superare il LIMITE. Propongo di iniziare con un set più piccolo e di aggiungere oggetti gradualmente'} \\ \tt \\ \tt PROPOSTA: {'B20', 'A11', 'C08', 'B93', 'A21'} 
	  } 
	   } 
	   } 
	  \\ 
 

    \theutterance \stepcounter{utterance}  

    & & \multicolumn{4}{p{0.6\linewidth}}{\cellcolor[rgb]{0.9,0.9,0.9}{%
	\makecell[{{p{\linewidth}}}]{% 
	  \tt {\tiny [A$\langle$GM]}  
	 RIFIUTO: {'A21', 'B93', 'A11', 'B33', 'C08', 'B39', 'B96', 'C85', 'A82', 'B31', 'B20', 'A60', 'B09', 'C81', 'B21'}\\ \tt \\ \tt ARGOMENTAZIONE: {'La proposta iniziale supera il LIMITE di 2307, quindi non è valida. Sono disposto a negoziare per trovare un set di oggetti che massimizzi l'importanza totale senza superare il LIMITE. Propongo di iniziare con un set più piccolo e di aggiungere oggetti gradualmente'} \\ \tt \\ \tt PROPOSTA: {'B20', 'A11', 'C08', 'B93', 'A21'} 
	  } 
	   } 
	   } 
	 & \\ 
 

    \theutterance \stepcounter{utterance}  

    & \multicolumn{4}{p{0.6\linewidth}}{\cellcolor[rgb]{0.8,1,0.9}{%
	\makecell[{{p{\linewidth}}}]{% 
	  \tt {\tiny [A$\rangle$GM]}  
	 RIFIUTO: {'B20', 'A11', 'C08', 'B93', 'A21'}\\ \tt \\ \tt ARGOMENTAZIONE: {'La tua proposta include oggetti con uno sforzo totale troppo elevato, considerando il LIMITE. Sostengo che alcuni oggetti abbiano un impatto maggiore sull'importanza totale e che dovremmo priorizzare quelli. Propongo di considerare oggetti con uno sforzo più basso, come "B33" o "B21", che potrebbero offrire un buon compromesso tra sforzo e importanza'}\\ \tt \\ \tt PROPOSTA: {'B33', 'B21', 'A82', 'B39'} 
	  } 
	   } 
	   } 
	 & & \\ 
 

    \theutterance \stepcounter{utterance}  

    & & \multicolumn{4}{p{0.6\linewidth}}{\cellcolor[rgb]{0.9,0.9,0.9}{%
	\makecell[{{p{\linewidth}}}]{% 
	  \tt {\tiny [GM$\rangle$B]}  
	 RIFIUTO: {'B20', 'A11', 'C08', 'B93', 'A21'}\\ \tt \\ \tt ARGOMENTAZIONE: {'La tua proposta include oggetti con uno sforzo totale troppo elevato, considerando il LIMITE. Sostengo che alcuni oggetti abbiano un impatto maggiore sull'importanza totale e che dovremmo priorizzare quelli. Propongo di considerare oggetti con uno sforzo più basso, come "B33" o "B21", che potrebbero offrire un buon compromesso tra sforzo e importanza'}\\ \tt \\ \tt PROPOSTA: {'B33', 'B21', 'A82', 'B39'} 
	  } 
	   } 
	   } 
	 & \\ 
 

    \theutterance \stepcounter{utterance}  

    & & & \multicolumn{4}{p{0.6\linewidth}}{\cellcolor[rgb]{1,0.85,0.72}{%
	\makecell[{{p{\linewidth}}}]{% 
	  \tt {\tiny [GM$\langle$B]}  
	 RIFIUTO: {'B33', 'B21', 'A82', 'B39'}\\ \tt \\ \tt ARGOMENTAZIONE: {'La tua nuova proposta è più ragionevole in termini di sforzo totale, ma ritengo che ci siano altre combinazioni di oggetti che potrebbero offrire un maggiore guadagno in termini di importanza totale. Sostengo che la selezione degli oggetti dovrebbe essere basata su una valutazione equilibrata tra lo sforzo richiesto e l'importanza percepita. Propongo di esplorare ulteriormente le opzioni, considerando anche oggetti con un sforzo moderato come "A60" o "B09", che potrebbero avere un impatto significativo sull'importanza totale'}\\ \tt \\ \tt PROPOSTA: {'B31', 'A60', 'B09', 'B33'} 
	  } 
	   } 
	   } 
	  \\ 
 

    \theutterance \stepcounter{utterance}  

    & & \multicolumn{4}{p{0.6\linewidth}}{\cellcolor[rgb]{0.9,0.9,0.9}{%
	\makecell[{{p{\linewidth}}}]{% 
	  \tt {\tiny [A$\langle$GM]}  
	 RIFIUTO: {'B33', 'B21', 'A82', 'B39'}\\ \tt \\ \tt ARGOMENTAZIONE: {'La tua nuova proposta è più ragionevole in termini di sforzo totale, ma ritengo che ci siano altre combinazioni di oggetti che potrebbero offrire un maggiore guadagno in termini di importanza totale. Sostengo che la selezione degli oggetti dovrebbe essere basata su una valutazione equilibrata tra lo sforzo richiesto e l'importanza percepita. Propongo di esplorare ulteriormente le opzioni, considerando anche oggetti con un sforzo moderato come "A60" o "B09", che potrebbero avere un impatto significativo sull'importanza totale'}\\ \tt \\ \tt PROPOSTA: {'B31', 'A60', 'B09', 'B33'} 
	  } 
	   } 
	   } 
	 & \\ 
 

    \theutterance \stepcounter{utterance}  

    & \multicolumn{4}{p{0.6\linewidth}}{\cellcolor[rgb]{0.8,1,0.9}{%
	\makecell[{{p{\linewidth}}}]{% 
	  \tt {\tiny [A$\rangle$GM]}  
	 ARGOMENTAZIONE: {'La tua proposta include una combinazione interessante di oggetti con sforzo moderato, che potrebbe offrire un buon equilibrio tra lo sforzo e l'importanza. Tuttavia, ritengo che "B31" possa essere sostituito con un oggetto che offre un rapporto sforzo-importanza più vantaggioso. Sostengo che la selezione degli oggetti debba essere basata su una valutazione olistica che consideri sia lo sforzo richiesto che l'importanza percepita. Propongo di valutare ulteriormente le opzioni per massimizzare l'importanza totale'}\\ \tt \\ \tt RIFIUTO: {'B31', 'A60', 'B09', 'B33'}\\ \tt \\ \tt PROPOSTA: {'A60', 'B09', 'B33', 'B21'} 
	  } 
	   } 
	   } 
	 & & \\ 
 

    \theutterance \stepcounter{utterance}  

    & & \multicolumn{4}{p{0.6\linewidth}}{\cellcolor[rgb]{0.9,0.9,0.9}{%
	\makecell[{{p{\linewidth}}}]{% 
	  \tt {\tiny [GM$\rangle$B]}  
	 ARGOMENTAZIONE: {'La tua proposta include una combinazione interessante di oggetti con sforzo moderato, che potrebbe offrire un buon equilibrio tra lo sforzo e l'importanza. Tuttavia, ritengo che "B31" possa essere sostituito con un oggetto che offre un rapporto sforzo-importanza più vantaggioso. Sostengo che la selezione degli oggetti debba essere basata su una valutazione olistica che consideri sia lo sforzo richiesto che l'importanza percepita. Propongo di valutare ulteriormente le opzioni per massimizzare l'importanza totale'}\\ \tt \\ \tt RIFIUTO: {'B31', 'A60', 'B09', 'B33'}\\ \tt \\ \tt PROPOSTA: {'A60', 'B09', 'B33', 'B21'} 
	  } 
	   } 
	   } 
	 & \\ 
 

    \theutterance \stepcounter{utterance}  

    & & & \multicolumn{4}{p{0.6\linewidth}}{\cellcolor[rgb]{1,0.85,0.72}{%
	\makecell[{{p{\linewidth}}}]{% 
	  \tt {\tiny [GM$\langle$B]}  
	 ARGOMENTAZIONE: {'La tua nuova proposta sembra essere una buona alternativa, poiché mantiene una combinazione di oggetti con sforzo moderato e aggiunge "B21", che potrebbe offrire un ulteriore guadagno in termini di importanza. Tuttavia, ritengo che valutare ulteriormente le opzioni per massimizzare l'importanza totale sia cruciale. Sostengo che ogni oggetto debba essere valutato in base al suo contributo unico all'importanza totale e che la selezione finale debba rappresentare il miglior compromesso possibile tra sforzo e importanza. Propongo di considerare se ci sono altre combinazioni di oggetti che potrebbero offrire un vantaggio maggiore in termini di importanza rispetto alla proposta attuale'}\\ \tt \\ \tt RIFIUTO: {'A60', 'B09', 'B33', 'B21'}\\ \tt \\ \tt PROPOSTA: {'A60', 'B09', 'B21', 'A82'} 
	  } 
	   } 
	   } 
	  \\ 
 

    \theutterance \stepcounter{utterance}  

    & & \multicolumn{4}{p{0.6\linewidth}}{\cellcolor[rgb]{0.9,0.9,0.9}{%
	\makecell[{{p{\linewidth}}}]{% 
	  \tt {\tiny [A$\langle$GM]}  
	 ARGOMENTAZIONE: {'La tua nuova proposta sembra essere una buona alternativa, poiché mantiene una combinazione di oggetti con sforzo moderato e aggiunge "B21", che potrebbe offrire un ulteriore guadagno in termini di importanza. Tuttavia, ritengo che valutare ulteriormente le opzioni per massimizzare l'importanza totale sia cruciale. Sostengo che ogni oggetto debba essere valutato in base al suo contributo unico all'importanza totale e che la selezione finale debba rappresentare il miglior compromesso possibile tra sforzo e importanza. Propongo di considerare se ci sono altre combinazioni di oggetti che potrebbero offrire un vantaggio maggiore in termini di importanza rispetto alla proposta attuale'}\\ \tt \\ \tt RIFIUTO: {'A60', 'B09', 'B33', 'B21'}\\ \tt \\ \tt PROPOSTA: {'A60', 'B09', 'B21', 'A82'} 
	  } 
	   } 
	   } 
	 & \\ 
 

    \theutterance \stepcounter{utterance}  

    & \multicolumn{4}{p{0.6\linewidth}}{\cellcolor[rgb]{0.8,1,0.9}{%
	\makecell[{{p{\linewidth}}}]{% 
	  \tt {\tiny [A$\rangle$GM]}  
	 ARGOMENTAZIONE: {'La tua ultima proposta introduce "A82", che potrebbe essere un oggetto chiave per massimizzare l'importanza totale. Sostengo che questa combinazione di oggetti sembra essere equilibrata e potrebbe offrire un buon compromesso tra sforzo e importanza. Tuttavia, è fondamentale considerare l'impatto complessivo di ogni oggetto sulla nostra strategia. Propongo di valutare se questa combinazione rappresenta il miglior set possibile per raggiungere l'obiettivo di massimizzare l'importanza totale senza superare il LIMITE'}\\ \tt \\ \tt ACCORDO: {'A60', 'B09', 'B21', 'A82'} 
	  } 
	   } 
	   } 
	 & & \\ 
 

    \theutterance \stepcounter{utterance}  

    & & & \multicolumn{2}{p{0.3\linewidth}}{\cellcolor[rgb]{0.95,0.95,0.95}{%
	\makecell[{{p{\linewidth}}}]{% 
	  \tt {\tiny [GM$|$GM]}  
	 game successful 
	  } 
	   } 
	   } 
	 & & \\ 
 

    \theutterance \stepcounter{utterance}  

    & & & \multicolumn{2}{p{0.3\linewidth}}{\cellcolor[rgb]{0.95,0.95,0.95}{%
	\makecell[{{p{\linewidth}}}]{% 
	  \tt {\tiny [GM$|$GM]}  
	 end game 
	  } 
	   } 
	   } 
	 & & \\ 
 

\end{supertabular}
}

\end{document}
